\section*{Author Section 3.2: proof of Theorem 1}
We show that if $ck<\zeta^5kQ^1_{p,\epsilon+12\zeta}(\widetilde f)/300-k\log\log(24/(5\zeta))$, the success of $\mathcal P$ is $(1-\epsilon)^{\Omega(\zeta^6k/\log|\mathcal Z|)}$.
\paragraph{Lemma 31.} Let $\delta=\zeta^6/1440000$, $\delta'=\zeta^6/(1440000\log|\mathcal Z|)$. Let $T_i$ indicate correctness of $\mathcal P$ on $f(X_i,Y_i)$. If its communication is $<\zeta^5kQ^1_{p,\epsilon+12\zeta}(\widetilde f)/60-k\log\log(24/(5\zeta))$, there exist $\lfloor\delta'k\rfloor$ coordinates $i_1,\ldots,i_{\lfloor\delta'k\rfloor}$ such that for all $1\le r\le\lfloor\delta'k\rfloor-1$ either
\[\Pr[\prod_{j=1}^rT_{i_j}=1]\le(1-\epsilon)^{\delta k},\]
or $\Pr[T_{i_{r+1}}=1\mid\prod_{j=1}^rT_{i_j}=1]\le1-\epsilon$.
Inductively suppose $C=\{i_1,\ldots,i_t\}$ has been found, $1\le t\le\lfloor\delta'k\rfloor$, with each conditional success at most $1-\epsilon$. Let $\mathcal E$ be $\prod_{i\in C}T_i=1$. If $\Pr[\mathcal E]\le(1-\epsilon)^{\delta k}$, done. Otherwise identify the next coordinate with conditional success at most $1-\epsilon$. The first coordinate is similar with empty conditioning; only the lower bound on conditioning probability is used.
Use $\ket\phi$, the state $\ket\rho$ conditioned on $\mathcal E$. For each $DG=dg$ define
\begin{multline*}\ket\phi_{X\widetilde XY\widetilde YABZ|dg}=\frac1{\sqrt{\gamma_{dg}}}\sum_{xy}\sqrt{P_{XY|dg}(xy)}\ket{xxyy}_{X\widetilde XY\widetilde Y}\\\otimes\sum_{z_C:(x_C,y_C,z_C)\in\widetilde f^t}\ket{z_C}_{Z_C}\ket{\widetilde\phi}_{ABZ_{\overline C}|xyz_C}.
\end{multline*}
The subnormalized state obeys $\|\ket{\widetilde\phi}_{ABZ_{\overline C}|xyz_C}\|_2^2=P_{Z_C|xy}(z_C)$. $\gamma_{dg}=\Pr[\mathcal E|dg]$ and $\sum_{dg}P_{DG}(dg)\gamma_{dg}=\Pr[\mathcal E]$. The distribution of $XYZ$ in this state is $P_{XYZ|\mathcal E,dg}$. Registers $DG$ are not explicitly considered: $\ket\phi_{dg}$ is notation without measuring such registers. Likewise $\ket\phi_{d_{-i}g_{-i}}$ conditions on $d_{-i}g_{-i}$ and changes only normalization. $\phi_{x_iy_id_{-i}g_{-i}}$ is obtained by measurement of the actual $X_iY_i$ registers. For $i\notin\overline C$, the authors use pure states $\ket\phi_{X_{\overline C}\widetilde X_{\overline C}Y_{\overline C}\widetilde Y_{\overline C}ABZ_{\overline C}|x_iy_ix_Cy_Cz_Cd_{-i}g_{-i}}$.
\paragraph{Lemma 32.} If $\Pr[\mathcal E]\ge(1-\epsilon)^{\delta k}$, there is $i\in\overline C$, $R_i=X_CY_CZ_CD_{-i}G_{-i}$, and for each $R_i=r_i$ a state $\ket{\phi'}_{X_{\overline C}\widetilde X_{\overline C}Y_{\overline C}\widetilde Y_{\overline C}ABZ_{\overline C}|y^*r_i}$ satisfying:
\begin{enumerate}\item[(i)] $\|P_{X_iY_iR_i|\mathcal E}-P_{X_iY_i}P_{R_i|\mathcal E,X_i}\|_1\le7\zeta/120$.
\item[(ii)] $\|P_{X_iY_iR_i|\mathcal E}-P_{X_iY_i}P_{R_i|\mathcal E,Y_i}\|_1\le7\zeta/120$.
\item[(iii)] There are projectors $\Pi_{x_ir_i}$ on $X_{\overline C}\widetilde X_{\overline C}A$ and unitaries $U_{y_ir_i}$ on $Y_{\overline C}\widetilde Y_{\overline C}BZ_{\overline C}$. Each projector succeeds on $\ket{\phi'}_{r_i}$ with probability $\alpha=2^{-c'}$, $c'\le60c/\zeta^5$, and
\begin{multline*}\mathbb E_{P_{X_iY_iR_i|\mathcal E}}\left[\left\|\frac1\alpha(\Pi_{x_ir_i}\otimes U_{y_ir_i})\proj{\phi'}_{y^*r_i}(\Pi_{x_ir_i}\otimes U_{y_ir_i}^\dagger)\right.\right.\\\left.\left.-\proj\phi_{x_iy_ir_i}\right\|_1\right]\le21\zeta.
\end{multline*}\end{enumerate}
Its complete author proof is in Section3.3 below.
\paragraph{Proof of Lemma 31.} Construct a protocol $\mathcal P'$ for $\widetilde f$ on distribution $P_{X_iY_i}=q$ by embedding inputs into the chosen state.
\begin{itemize}\item Alice/Bob have $r$ distributed as required by Fact12 as shared randomness, and $2^{60c/\zeta^5}\log(24/(5\zeta))$ copies of $\ket{\phi'}_{y^*r_i}$ as entanglement, Alice holding $X_{\overline C}\widetilde X_{\overline C}A$ and Bob the remaining registers.
\item On input $(x_i,y_i)$, Fact12 and Lemma32(i),(ii) give $R_i^A,R_i^B$ with
\[\|P_{X_iY_iR_i^AR_i^B}-P_{X_iY_iR_iR_i|\mathcal E}\|_1\le7\zeta/24,\]
where $R_iR_i$ are perfectly correlated copies.
\item Alice measures $\{\Pi_{x_ir_i^A},1-\Pi_{x_ir_i^A}\}$ on each copy. If none succeeds, abort. Otherwise send a $(60c/\zeta^5+\log\log(24/(5\zeta)))$-bit message indexing a successful copy.
\item Bob applies $U_{y_ir_i^B}$ on that copy and measures $Z_i$ for his output.\end{itemize}
By construction,
\[\mathbb E_{P_{X_iY_iR_i|\mathcal E}}\Pr[Z_i\text{ measured on }\ket\phi_{x_iy_ir_i}\in\widetilde f(x_i,y_i)]=\Pr[T_i=1|\mathcal E].\]
First suppose the inputs and $R_i^A,R_i^B$ follow exactly $P_{X_iY_iR_iR_i|\mathcal E}$, writing both copies $r_i$. The abort probability is
\[(1-2^{-c'})^{2^{60c/\zeta^5}\log(24/(5\zeta))}\le5\zeta/24.\]
If not aborting, the state after Bob's unitary is $\alpha^{-1/2}(\Pi_{x_ir_i}\otimes U_{y_ir_i})\ket{\phi'}_{y^*r_i}$. Lemma32(iii) gives expected correctness at least $\Pr[T_i=1|\mathcal E]-21\zeta/2$. Including abort, at least $\Pr[T_i=1|\mathcal E]-21\zeta/2-5\zeta/24$. For the actual $R_i^A,R_i^B$ distribution, success is at least
\[\Pr[T_i=1|\mathcal E]-21\zeta/2-5\zeta/24-7\zeta/24\ge\Pr[T_i=1|\mathcal E]-11\zeta.\]
Since $\|q-p\|_1\le2\zeta$, running this protocol on $p$ gives success $\ge\Pr[T_i=1|\mathcal E]-12\zeta$. Communication is $60c/\zeta^5+\log\log(24/(5\zeta))<Q^1_{p,\epsilon+12\zeta}(\widetilde f)$. If $\Pr[T_i=1|\mathcal E]\ge1-\epsilon$, error is $\le\epsilon+12\zeta$, a contradiction. Hence $\Pr[T_i=1|\mathcal E]\le1-\epsilon$. Take $i_{t+1}=i$.\qed
